\documentclass[10pt,a4paper]{amsart}
\usepackage[margin=19mm]{geometry}
\usepackage{amsmath,amssymb,mathtools}
\usepackage[scr=rsfs,cal=euler]{mathalfa}
\usepackage{xfrac}
\usepackage[unicode,hidelinks,bookmarks=false]{hyperref}
\newcommand{\ie}{i.e.\ }
\newcommand{\cf}{cf.\ }
\newcommand{\resp}{respectively\ }
\newcommand{\Subsection}{\subsection}
\providecommand{\nobreakdash}{\nobreak\hbox{-}}
\setlength{\emergencystretch}{2em}
\multlinegap0pt
\usepackage{bm}
\usepackage{bbm}
\usepackage{dsfont}
\usepackage{xspace}
\usepackage{cancel}
\usepackage{mathtools}
\usepackage{amsthm}
\makeatletter
\@ifundefined{lemma}{\newtheorem{lemma}{Lemma}}{}
\@ifundefined{theorem}{\newtheorem{theorem}{Theorem}}{}
\makeatother
\title[Dvorak-Dell-Grohe-Rattan theorem via an asymptotic argument]{Dvorak-Dell-Grohe-Rattan theorem via an asymptotic argument}
\author{Alexander Kozachinskiy}
\date{Source: arXiv:2507.14669v1 (CC BY 4.0)}
\begin{document}
\maketitle

\noindent\emph{Source and licence.} Dvorak-Dell-Grohe-Rattan theorem via an asymptotic argument. Version arXiv:2507.14669v1 (CC BY 4.0), distributed under the Creative Commons Attribution 4.0 International licence, as recorded in the arXiv licence metadata for this exact version and shown on the abstract page https://arxiv.org/abs/2507.14669v1. Full rights notice: licenses/arxiv-2507.14669-CC-BY-4.0.txt.\par\smallskip

\noindent\textit{Editorial scope.} Counted unit: the Theorem whose source label is \texttt{None} in the article (source0.tex lines 73--75, source label \texttt{None}), delivered together with the article's own complete original proof of it (source0.tex lines 76--131, one \texttt{proof} environment of 56 lines). No mathematics is written, repaired or re-derived: statement and proof are the article's own text line for line. Required context retained uncounted: none. Every definition, display and label of those lines is retained unchanged. Omitted from this package and not redistributed: 1--72, 132--159. Nothing inside a retained excerpt is omitted or altered: every retained body is delivered exactly as the article's lines are written, so no omission receipt of any kind accompanies this package. Citations: the retained excerpts carry no citation command at all, so no bibliography entry of the article is transcribed and no reference is redistributed. Reference adaptation: none was needed and none was made; every reference inside the delivered unit resolves inside it, so no unresolved reference or citation is produced. Printed slips kept literal: no printed word, symbol, hypothesis or number of the counted statement or of its proof was altered. Layout and class: the article's own class is \texttt{article} with packages including graphicx, hyperref, amsmath, amsthm, amssymb, tikz, tcolorbox, dashrule, xcolor, wrapfig, fullpage, algorithm2e; the wrapper is the lane's pinned \texttt{amsart} editorial wrapper at 10pt on A4, which restates the theorem-like environments the retained text needs and adds the article's own macro definitions that the retained text needs (none), with extra wrapper packages (bm, bbm, dsfont, xspace, cancel, mathtools). The delivered numbering is the unit's own. Assets: no retained span contains a figure environment, a graphics-inclusion command, a PDF-page inclusion, a table or a TikZ picture; no image file is copied into this package, the package declares no asset dependency and no image file is redistributed with it. Licence: version arXiv:2507.14669v1 of this article is distributed under the Creative Commons Attribution 4.0 International licence, as recorded in the arXiv licence metadata for this exact version and shown on the abstract page https://arxiv.org/abs/2507.14669v1. A full rights notice is delivered at licenses/arxiv-2507.14669-CC-BY-4.0.txt. Characters: every retained excerpt is pure ASCII, so the delivered engine, XeLaTeX with its default Latin Modern text font, reports no missing character.
\begin{theorem}
    If two graphs $G_1$ and $G_2$ are distinguished by the WL test, then there is a tree $T$ that has a different number of homomorphisms to $G_1$ and to $G_2$.
\end{theorem}

\begin{proof}




    For the proof, we define, for any $k$, a linear order on $k$-level WL-labels. For $k = 0$, there is just one possible label, so there is nothing to define. For $k = 1$, each label is identified with a natural number (degree of a node in a graph). The larger this number is, the larger we put the label in our order.

    More generally, assume we have already defined the order on $k$-level labels. We extend it to $(k+1)$-level labels ``lexicographically''. Namely, $(k+1)$-level labels are multisets of $k$-level labels. Now, we take two $(k+1)$-labels $\ell_1$ and $\ell_2$ that we want to ``compare''. There is already an order on their elements (these elements are $k$-level labels). We take the largest element that appears a different number of times in the multisets $\ell_1$ and $\ell_2$. The multiset that has this element more times is defined to be the larger one. This is clearly a linear order.


\begin{lemma}
\label{lem_dom}
    For any finite set $S$ of $WL$-labels that does not have labels of isolated nodes, and for any $k$, there exists  a sequence of $k$-depth rooted trees $\{T^k_n\}_{n\to\infty}$ such that for any two $k$-level labels $\ell_1, \ell_2\in S$, we have:
     \begin{equation}
     \label{eq_order}
         \ell_1 < \ell_2 \implies h(T^k_n, \ell_1) = o(h(T^k_n, \ell_2)) \text{ as } n\to\infty.\
     \end{equation}
\end{lemma}

%\begin{remark}
 %   Since $S$ does not have labels of isolated nodes, we have $h(T, \ell) \ge 1$ for any tree $T$ and for any $\ell \in L$. Hence, the conditions  $h(T^k_n, \ell_1) = o(h(T^k_n, \ell_2)) \text{ as } n\to\infty$ are well-defined, because the fraction is $h(T^k_n, \ell_1)/h(T^k_n, \ell_2)$ is well-defined.
%\end{remark}


    
    Let us first deduce the theorem from the lemma. Take any two graphs $G_1$ and $G_2$ that are distinguished after the $k$-th iteration of the WL-test.  Assume first, for simplicity, that $G_1$ and $G_2$ do not have isolated nodes.
    We apply the lemma to the set $L$of $k$-level WL labels that appear in $G_1$ and $G_2$. We will show that one of the trees $T^k_n$ distinguishes $G_1$ and $G_2$ by the number of homomorphisms from it (as from a non-rooted tree). Indeed, the number of homomorphisms from $T^k_n$ to a graph $G$ can be counted as the sum over vertices $v$ of $G$, of the number of homomorphisms from $T^k_n$ to $G$ that map the root to $v$. Hence, it is enough to show that for a large enough $n$, we have:
    \begin{equation}
    \label{eq_dif}
    \sum\limits_{v\in V(G_1)} h(T^k_n,\ell^k(v)) \neq \sum\limits_{v\in V(G_2)} h(T^k_n,\ell^k(v)). 
    \end{equation}
    Graphs $G_1$ and $G_2$ have different multisets of $k$-level labels. Consider the largest (in our ``lexicographic'' order) $k$-level label $\ell$ that appears a different number of times in them. Let us say, it appears more in $G_1$. If we consider the difference between the left-hand side and the right-hand side in \eqref{eq_dif}, all terms corresponding to labels that are larger than $\ell$ will cancel out. In turn, $h(T_n^k, \ell)$ will be with a positive coefficient in this difference, and $h(T_n^k,\ell)$ is itself strictly positive because $\ell$ is for a non-isolated node. Hence, for large enough $n$, the difference will be strictly positive. Indeed, terms $h(T_n^k, \ell_1)$ for $\ell_1 < \ell$ can have negative coefficients, but because these terms are dominated by  $h(T_n^k, \ell)$, this will not prevent the difference from being positive in the limit.

    What if $G_1$ and $G_2$ have isolated nodes? Then in \eqref{eq_dif}, some terms are 0s, corresponding to isolated nodes. If the multisets of non-isolated nodes are different, the same argument as before applies. If they are the same, then the graphs have to differ just in the number of nodes, in which case we can distinguish them by the number of homomorphisms from a single isolated node.


    \begin{proof}[Proof of Lemma \ref{lem_dom}]
        Let us first establish the lemma for $k = 1$. We will let $T^1_n$ be just a star with $n$ children. There are $d^n$ homomorphisms from $T^1_n$ to a node of degree $d$ that map the root to this node (for every child of the root, there are $d$ options). Since $d_1^n = o(d_2^n)$ for $1\le d_1 < d_2$, the induction base follows.

We now perform the induction step. Assume that the statement is proved for $k$. First, let us for every $n$ define a tree $H_n$ of depth $k + 1$  where the root has one child whose subtree is $T_n^k$. Let $L$ be a $(k + 1)$-level WL-label, i.e., a multiset of some $k$-level WL-labels $\ell$. We have:
\[h(H_n, L) =  \sum\limits_{\ell\in L} h(T_n^k, \ell).\]
Indeed, if $v$ is a node with label $L$, then we have to map the unique child of the root of $H_n$ to some neighbor of $v$. Upon choosing this neighbor, we have to map $T_n^k$ to this neighbor, and the number of ways to do that is $h(T_n^k, \ell)$, where $\ell$ is the $k$-level WL-label of this neighbor. And by definition, the multiset of $k$-level WL-labels of neighbors of $v$ is $L$. 


 Now, there exists $m$ such that $h(H_m, L_1) < h(H_m, L_2)$ for all $L_1 < L_2 \in S$. Indeed, based on the condition \eqref{eq_order} for level $k$, we have for any two fixed $L_1, L_2 \in S$ with $L_1 < L_2$ 
 \[h(H_n, L_1) < h(H_n, L_2)\]
 for all large enough $n$ (we consider the largest $k$-level label $\ell$ that appears a different number of times in $L_1$ and $L_2$, there are fewer of them in $L_1$, so the difference  $h(H_n, L_2) - h(H_n, L_1)$ includes this term with a positive coefficient, this term is itself positive because of the assumption about not having isolated nodes, and all other terms are dominated by it). There are finitely many labels in $S$, so it remains to take $m$ large enough so that this inequality holds for all pairs in question.


 What remains to do is to define $T_n^{(k+1)}$ as a depth-$(k+1)$ tree where the root has $n$ children, each followed by $T_m^k$ (or, in other words, $n$ copies of $H_m$, glued by the root). Observe that 
 \[h(T_n^{k+1}, L) = h(H_m, L)^n.\]
Taking any two $L_1, L_2 \in S$ with $L_1 < L_2$, since  $h(H_m, L_1) < h(H_m, L_2)$, we obtain:
\[h(T_n^{k+1}, L_1) = h(H_m, L_1)^n = o( h(H_m, L_2)^n) = o(h(T_n^{k+1}, L_2)),\]
as required.
    \end{proof}
\end{proof}

\end{document}
